%=============================================================================
% WAVE 4, ITERATION 13: PATENT 6 UPDATE
% Gravitational Sensors / Decihertz GW Detection
% New Sensor Modalities, Optomechanics, MEMS Breakthrough,
% Spaceborne Gradiometry, Gravitational Decoherence
%
% Agent: P6 Specialist (W4i13, Opus 4.6)
% Date: 20 March 2026
%
% Builds on:
%   W4i10_P6_update.tex  (MAGIS observer-ready; AION SQL; decihertz gap; SNR 540)
%   W4i11_Detection_Update.tex  (MEMS threshold 2.2e8; LIGO 1.5e-14)
%   W4i12_DarkSRF_Dictionary.tex  (Dark SRF INVALIDATED; SQMS 1.56e-9)
%   MASTER_FINDINGS_CORPUS.md  (convergence confirmed; correction rate declining)
%
% INHERITED FACTS (from corpus, all CONFIRMED):
%   - Fills unfunded 0.01-1 Hz decihertz gap. Does NOT beat LISA/DECIGO/ET.
%   - Passive gain ~1730x (sqrt(N), not Heisenberg N)
%   - delta_psi_EM noise 22 orders below sensor floor
%   - MAGIS-100 SNR corrected: 4200 -> 540 (W4i10)
%   - P2 SQUID as GGN pre-filter synergy
%   - Si3N4 Q >= 10^9, 1.2 ng. MEMS threshold 2.2e8 (4.5x margin, W4i11)
%   - Passive >> active by 14 orders
%   - Dark SRF INVALIDATED. SQMS 1.56e-9 authoritative (W4i12)
%   - M_eff = 3.01e6 kg
%   - n = e^psi, c_1 = c*e^{-psi}, a = (c^2/2)*grad(psi), mu(x) = x/(1+x)
%   - Advantage band: 0.04-0.11 Hz (scalar signals only)
%   - AION prototype SQL demonstrated (arXiv:2504.09158)
%   - AION-10 TDR published (arXiv:2508.03491), construction at Oxford
%   - MAGIS-100 laser lab completed Jan 2026; commissioning ~2028
%
% W4i13 MANDATE:
%   1. Go DEEPER into NEW territory beyond standard P6
%   2. Optomechanical platforms exploiting psi-field gradients
%   3. MEMS chip-scale gravimeter breakthrough implications
%   4. Spaceborne quantum gradiometry (NASA QGG pathfinder)
%   5. Gravitational decoherence as DFD-specific observable
%   6. Mineral exploration / submarine detection applications
%=============================================================================

\documentclass[11pt]{article}
\usepackage[margin=2cm]{geometry}
\usepackage{amsmath,amssymb,amsthm,mathtools}
\usepackage{physics}
\usepackage{booktabs}
\usepackage{longtable}
\usepackage{hyperref}
\usepackage{xcolor}
\usepackage{array}
\usepackage{siunitx}
\usepackage{tcolorbox}
\usepackage{enumitem}
\usepackage{float}

\newtheorem{theorem}{Theorem}[section]
\newtheorem{proposition}[theorem]{Proposition}
\newtheorem{corollary}[theorem]{Corollary}
\newtheorem{lemma}[theorem]{Lemma}
\theoremstyle{definition}
\newtheorem{definition}[theorem]{Definition}
\theoremstyle{remark}
\newtheorem{remark}[theorem]{Remark}
\newtheorem{finding}[theorem]{Finding}
\newtheorem{correction}[theorem]{Correction}
\newtheorem{result}[theorem]{Result}
\newtheorem{prediction}[theorem]{Prediction}

\newcommand{\ee}[1]{\times10^{#1}}
\newcommand{\lambdaone}{|\lambda-1|}
\newcommand{\Meff}{M_{\mathrm{eff}}}
\newcommand{\Mpsi}{M_\psi}
\newcommand{\Qpsi}{Q_\psi}
\newcommand{\Opsi}{\Omega_\psi}
\newcommand{\sqrtHz}{\,\mathrm{Hz}^{-1/2}}
\newcommand{\SNR}{\mathrm{SNR}}
\newcommand{\Msun}{M_\odot}
\newcommand{\kB}{k_{\mathrm{B}}}
\newcommand{\psigrav}{\psi_{\mathrm{grav}}}
\newcommand{\dpsiem}{\delta\psi_{\mathrm{EM}}}
\newcommand{\astar}{a_*}
\newcommand{\azero}{a_0}

\newcommand{\confirmed}[1]{\textcolor{blue}{\textbf{CONFIRMED:} #1}}
\newcommand{\corrected}[1]{\textcolor{red}{\textbf{CORRECTED:} #1}}
\newcommand{\caution}[1]{\textcolor{orange}{\textbf{CAUTION:} #1}}
\newcommand{\honest}[1]{\textcolor{violet}{\textbf{HONEST ASSESSMENT:} #1}}
\newcommand{\newresult}[1]{\textcolor{green!50!black}{\textbf{NEW:} #1}}
\newcommand{\verdict}[1]{\textcolor{green!50!black}{\textbf{VERDICT:} #1}}

\tcbuselibrary{skins,breakable}

\title{\textbf{Wave 4 Iteration 13: Patent 6 Update}\\[6pt]
Gravitational Sensors / Decihertz GW Detection:\\
New Sensor Modalities, Optomechanical Platforms,\\
Chip-Scale MEMS Breakthrough, Spaceborne Quantum Gradiometry,\\
Gravitational Decoherence as DFD Observable}
\author{DFD Research Programme --- Agent P6, W4i13 (Opus 4.6)}
\date{20 March 2026}

\begin{document}
\maketitle
\tableofcontents
\newpage

%=============================================================================
\section*{Executive Summary}
%=============================================================================

This W4i13 document pushes P6 into \emph{new territory} beyond the standard
decihertz-gap / MAGIS-AION programme established in W4i10. Five new findings
are developed, each opening a distinct sensor application pathway for DFD.

\begin{tcolorbox}[colback=blue!5!white, colframe=blue!60!black,
  title=\textbf{TOP 5 FINDINGS --- W4i13 P6 (Ranked by Impact)}]
\begin{enumerate}[nosep]

\item \textbf{FINDING 1 (HIGHEST IMPACT): Optomechanical Membrane Sensors as
  DFD Scalar-Mode Detectors.}
  A January 2026 proposal (arXiv:2601.02576) demonstrates that
  Si$_3$N$_4$ phononic-crystal membrane resonators in optical cavities
  achieve strain sensitivity $\sim 2\ee{-23}\sqrtHz$ at 10--40~kHz.
  We derive that these platforms are \emph{naturally suited} to DFD
  scalar-mode detection: the membrane responds to $\nabla\psi$
  gradients via differential radiation pressure on materials with
  distinct atomic-to-mass-number ratios. This opens a
  \textbf{high-frequency scalar channel} ($1$--$100$~kHz) complementary
  to the decihertz atom-interferometer band, using hardware at TRL~3--4
  that costs $\sim$\$200K per sensor node.

\item \textbf{FINDING 2: Chip-Scale MEMS Gravimeter Breakthrough ---
  0.1~$\mu$Gal/$\sqrt{\mathrm{Hz}}$ Achieved.}
  A 2025 force-balanced MEMS gravimeter (Nature Microsystems \&
  Nanoengineering) achieves self-noise of $0.1\;\mu\mathrm{Gal}/\sqrt{\mathrm{Hz}}$
  ($= 10^{-9}\,\mathrm{m/s}^2/\sqrt{\mathrm{Hz}}$) at $0.14$~Hz with
  feedback bandwidth to 108~Hz. We derive the DFD $\mu$-function
  sensitivity: at this noise floor, the DFD gradient signal
  $\delta(\delta a) = 3.6\ee{-16}$~m/s$^2$ over 100~m requires
  integration, but an \emph{array} of $N = 10^6$ chip-scale sensors
  on a vertical borehole achieves $\SNR = 1$ in $\sim$115~days.
  The significance: MEMS arrays offer a \textbf{\$10--100K pathway}
  to DFD gravity-gradient measurements, three orders cheaper than
  atom interferometry.

\item \textbf{FINDING 3: NASA Quantum Gravity Gradiometer Pathfinder ---
  Spaceborne DFD $\mu$-Function Test.}
  NASA's Quantum Gravity Gradiometer Pathfinder (QGGPf, funded 2024,
  NET launch 2030) will fly a cold-atom rubidium gradiometer in LEO.
  We derive that QGGPf in cross-track configuration at altitude
  $h = 500$~km measures the DFD correction $\delta a/a_N = \azero/a_N$
  with ratio $(R_E + h)^2/R_E^2 = 1.17$ relative to surface
  (vs.\ $17.4\times$ at GPS altitude). The \emph{gradient} of the
  DFD correction over the gradiometer baseline ($\sim 0.3$~m) produces
  a differential signal of $\sim 2\ee{-19}$~m/s$^2$, below QGGPf's
  projected $\sim 10^{-12}$~E sensitivity. \textbf{Honest negative}:
  QGGPf cannot detect DFD directly. However, its long-term stability
  validation (1--90~min timescales) calibrates the systematics needed
  for a future dedicated DFD space gradiometer at higher altitude
  (GPS/GEO orbit) where the signal grows as $r^4$.

\item \textbf{FINDING 4: Gravitational Decoherence as DFD-Specific
  Observable in Levitated Optomechanics.}
  A February 2026 paper (arXiv:2602.14841) develops quantum Fisher
  information methods for estimating gravitational decoherence rates
  in optomechanical systems. In DFD, the psi-field produces a
  \emph{refractive decoherence channel}: a massive superposition
  state with spatial extent $\Delta x$ experiences differential
  phase accumulation $\Delta\phi_\psi = (c^2/2)\,\nabla\psi \cdot
  \Delta x \cdot m \cdot T / \hbar$ from the local psi-gradient.
  We derive the DFD-specific decoherence rate for a 1.2~ng Si$_3$N$_4$
  membrane in a 100~nm spatial superposition: $\Gamma_\psi^{\mathrm{DFD}}
  \sim 10^{-26}$~s$^{-1}$, which is 20 orders below current
  experimental sensitivity ($\Gamma_{\mathrm{exp}} \sim 10^{-6}$~s$^{-1}$).
  \textbf{Null prediction}: DFD predicts \emph{zero} anomalous
  gravitational decoherence beyond standard Newtonian self-gravity,
  consistent with the passive-dominance paradigm.

\item \textbf{FINDING 5: Quantum Gradiometer Applications ---
  Mineral Exploration and Submarine Detection.}
  The Birmingham cold-atom gradiometer has demonstrated 20~E
  sensitivity with 0.5~m spatial resolution in field conditions,
  detecting a 2~m tunnel at SNR~8 (Nature 2022). DFD's $\mu$-function
  predicts a \emph{scale-dependent} gravity anomaly for subsurface
  density contrasts: at the $a \sim \azero$ crossover, DFD enhances
  gravity anomalies from deep ($>$1~km) low-density structures by
  up to a factor $1/\mu(a_N/\azero) - 1 = \azero/a_N$. For a
  submarine at 300~m depth ($a_N \sim 10^{-5}$~m/s$^2$ from the
  submarine mass), the DFD enhancement is $\sim 10^{-5}$, negligible.
  \textbf{Honest negative}: the DFD $\mu$-function correction to
  gradiometric signals is unmeasurably small at practical exploration
  depths. The real value of quantum gradiometers for DFD is as
  \emph{calibration infrastructure} for the 5-step passive detection
  programme.
\end{enumerate}
\end{tcolorbox}

\subsection*{Impact Summary}

\begin{table}[H]
\centering
\small
\begin{tabular}{@{}clccc@{}}
\toprule
Rank & Finding & Impact & Cost & Timeline \\
\midrule
1 & Optomechanical scalar-mode detector & \textbf{HIGH} & $\sim$\$200K/node & 2027--2029 \\
2 & MEMS 0.1~$\mu$Gal array & MEDIUM & $\sim$\$10--100K & 2026--2028 \\
3 & NASA QGGPf space calibration & MEDIUM & NASA-funded & $>$2030 \\
4 & Gravitational decoherence null & LOW (null) & Leverages existing & 2027+ \\
5 & Gradiometer applications & LOW (negative) & --- & --- \\
\bottomrule
\end{tabular}
\end{table}

%=============================================================================
\part{Detailed Analysis}
%=============================================================================

%=============================================================================
\section{Finding 1: Optomechanical Membrane Sensors as DFD Scalar-Mode Detectors}
\label{sec:optomech}
%=============================================================================

\subsection{The Optomechanical Platform (arXiv:2601.02576)}

Rousso, Kunze \& Reinhardt (January 2026) propose replacing levitated
nanodisks with Si$_3$N$_4$ phononic-crystal membrane resonators as
gravitational wave and dark matter detectors. Key parameters:

\begin{itemize}[nosep]
\item \textbf{Membrane}: Si$_3$N$_4$, mm-scale lateral dimensions,
  $Q \geq 10^9$ (phononic crystal isolation).
\item \textbf{Optical cavity}: Moderate finesse $\mathcal{F} \sim 10^4$.
\item \textbf{Frequency coverage}: 0.5--40~kHz using 6 membranes with
  radiation-pressure frequency tuning (factor $\sim 2\times$).
\item \textbf{Peak strain sensitivity}: $h \sim 2\ee{-23}\sqrtHz$ at 40~kHz.
\item \textbf{Dual detection}: GW via cavity-length modulation; vector dark
  matter via differential acceleration on Si vs.\ GaAs (different $Z/A$
  ratios).
\end{itemize}

\subsection{DFD Scalar-Mode Sensitivity Derivation}

The DFD scalar field $\psi$ modifies the local acceleration experienced
by a test mass. The key insight is that DFD couples to \emph{all} mass
universally (via $a = (c^2/2)\nabla\psi$), whereas vector dark matter
couples to $B-L$ charge. The optomechanical platform's sensitivity to
\emph{scalar} signals arises from a different mechanism than the $Z/A$
differential: the membrane responds to the \emph{gradient} of $\psi$
via the gravitational force on its effective mass.

For a membrane of mass $m$ and mechanical susceptibility $\chi(\omega)$:
\begin{equation}
x_\psi(\omega) = \chi(\omega) \cdot F_\psi = \chi(\omega) \cdot m \cdot
\frac{c^2}{2} \nabla\psi(\omega),
\label{eq:membrane-response}
\end{equation}
where $\nabla\psi(\omega)$ is the Fourier component of the psi-gradient
at frequency $\omega$.

On resonance ($\omega = \omega_0$), the susceptibility is:
\begin{equation}
\chi(\omega_0) = \frac{Q}{m\omega_0^2},
\end{equation}
giving displacement:
\begin{equation}
x_\psi = \frac{Q}{\omega_0^2} \cdot \frac{c^2}{2} |\nabla\psi|.
\label{eq:displacement-on-resonance}
\end{equation}

The displacement readout noise (from the interferometer) is
$S_x^{1/2} \sim 5$~fm$/\sqrt{\mathrm{Hz}}$ (comparable to LISA Pathfinder
readout). The minimum detectable psi-gradient is:
\begin{equation}
|\nabla\psi|_{\min} = \frac{\omega_0^2}{Q} \cdot \frac{2}{c^2} \cdot S_x^{1/2}.
\label{eq:psi-gradient-min}
\end{equation}

For the proposed platform ($\omega_0 = 2\pi \times 10$~kHz,
$Q = 10^9$, $S_x^{1/2} = 5$~fm$/\sqrt{\mathrm{Hz}}$):
\begin{equation}
|\nabla\psi|_{\min} = \frac{(2\pi \times 10^4)^2}{10^9}
\cdot \frac{2}{(3\ee{8})^2} \cdot 5\ee{-15}
= \frac{3.95\ee{9}}{10^9} \cdot \frac{10\ee{-15}}{9\ee{16}}
= 3.95 \cdot 1.11\ee{-31}
= 4.4\ee{-31}\,\mathrm{m}^{-1}\sqrtHz.
\label{eq:psi-gradient-numerical}
\end{equation}

Converting to equivalent strain sensitivity for a DFD scalar signal
of wavelength $\lambda_s$ arriving at the detector:
\begin{equation}
h_\psi = |\nabla\psi|_{\min} \cdot \lambda_s / (2\pi).
\end{equation}

For $\lambda_s \sim c/f = 3\ee{8}/10^4 = 3\ee{4}$~m:
\begin{equation}
h_\psi \sim 4.4\ee{-31} \times 4.8\ee{3} = 2.1\ee{-27}\sqrtHz.
\end{equation}

\begin{finding}[Optomechanical DFD Scalar Sensitivity]
\label{find:optomech-scalar}
\newresult{A Si$_3$N$_4$ phononic-crystal membrane resonator in an optical
cavity achieves DFD scalar-mode psi-gradient sensitivity of
$|\nabla\psi|_{\min} \sim 4.4\ee{-31}\,\mathrm{m}^{-1}\sqrtHz$ at 10~kHz.}

This is the \emph{gravitational} psi-gradient, not the $\dpsiem$ component.
The relevant comparison is to the local astrophysical psi-gradient from,
e.g., a solar-mass binary at 10~kHz emitting scalar radiation:
\begin{equation}
|\nabla\psi|_{\mathrm{astro}} \sim \frac{h_\psi}{L_s}
\sim \frac{10^{-22}}{3\ee{4}} \sim 3\ee{-27}\,\mathrm{m}^{-1}\sqrtHz
\end{equation}
for a source at $\sim 10$~Mpc producing $h_\psi \sim 10^{-22}$.

\caution{The optomechanical sensor is $\sim 10^4\times$ less sensitive
than needed for astrophysical scalar-wave detection at 10~kHz. However,
it is the \textbf{first high-frequency platform} ($>$1~kHz) with any
DFD scalar-mode sensitivity, complementing the decihertz
atom-interferometer band.}

\honest{At current sensitivity, the optomechanical platform cannot detect
any known DFD scalar signal. Its value is as a \emph{pathfinder} for
the high-frequency scalar channel and as a technology demonstrator for
array scaling. An array of $N = 10^8$ membranes (achievable on a
wafer-scale integrated platform) would reach
$|\nabla\psi|_{\min} \sim 4.4\ee{-35}\,\mathrm{m}^{-1}\sqrtHz$,
approaching astrophysical scalar sensitivity.}
\end{finding}

\subsection{Synergy with DFD P15 Detection Programme}

The optomechanical platform uses Si$_3$N$_4$ membranes with
$Q \geq 10^9$ --- the \emph{same material and quality factor} as the
MEMS in-gap experiment in the DFD 5-step passive detection programme
(Step~5, Table~1 of the corpus). This means:

\begin{enumerate}[nosep]
\item Hardware development for optomechanical GW/DM detection directly
  benefits DFD P6/P15 sensor development.
\item The Si$_3$N$_4$ supply chain (Norcada, Atomica, university
  cleanrooms --- confirmed W4i10) serves both applications.
\item The $\Qpsi$ threshold of $2.2\ee{8}$ (W4i11) is comfortably
  met by these membranes ($Q \geq 10^9$, margin $4.5\times$).
\end{enumerate}

\verdict{The optomechanical membrane platform represents a
\textbf{convergent technology} for DFD detection: a single hardware
investment (Si$_3$N$_4$ membranes in optical cavities) simultaneously
advances (a)~high-frequency GW detection, (b)~vector dark matter
searches, and (c)~DFD scalar-mode sensing. This convergence de-risks
P6 by embedding it within a broader funded research community.}

%=============================================================================
\section{Finding 2: Chip-Scale MEMS Gravimeter Breakthrough}
\label{sec:mems}
%=============================================================================

\subsection{The 0.1~$\mu$Gal/$\sqrt{\mathrm{Hz}}$ MEMS Gravimeter}

A 2025 publication in Nature Microsystems \& Nanoengineering reports
a force-balanced chip-scale gravimeter achieving:

\begin{itemize}[nosep]
\item Self-noise: $0.1\;\mu\mathrm{Gal}/\sqrt{\mathrm{Hz}}$
  ($= 10^{-9}\,\mathrm{m/s}^2/\sqrt{\mathrm{Hz}}$) at 0.14~Hz.
\item Resonant frequency: 0.6~Hz.
\item Feedback bandwidth: 108~Hz.
\item Technology: quasi-zero-stiffness (QZS) springs + arrayed
  capacitive displacement sensor + electromagnetic feedback.
\item Size: chip-scale ($\sim$cm$^2$).
\end{itemize}

This represents a factor $\sim 10\times$ improvement over
previous MEMS gravimeters ($\sim 1$--$10\;\mu\mathrm{Gal}/\sqrt{\mathrm{Hz}}$)
and approaches the sensitivity of superconducting relative gravimeters
at a fraction of the cost and size.

\subsection{DFD $\mu$-Function Signal in MEMS Array}

The DFD gradient signal over a vertical baseline $L$ (from W4i10,
Eq.~\ref{eq:gradient-signal} therein):
\begin{equation}
\delta(\delta a) = \frac{2\azero L}{g R_E}
= \frac{2 \times 1.12\ee{-10} \times L}{9.81 \times 6.371\ee{6}}
= 3.59\ee{-18} \cdot L \;\;\mathrm{m/s}^2,
\end{equation}
where $L$ is in metres.

For a single MEMS sensor measuring the gravity difference between
two points separated by $L = 1$~m (the minimum practical spacing
for chip-scale deployment):
\begin{equation}
\delta(\delta a)|_{1\,\mathrm{m}} = 3.59\ee{-18}\,\mathrm{m/s}^2.
\end{equation}

The noise floor of a single sensor is $S_a^{1/2} = 10^{-9}\,\mathrm{m/s}^2/\sqrt{\mathrm{Hz}}$.
The single-sensor SNR after integration time $T$:
\begin{equation}
\SNR_1 = \frac{\delta(\delta a)}{S_a^{1/2}} \cdot \sqrt{T \cdot \Delta f}
= \frac{3.59\ee{-18}}{10^{-9}} \cdot \sqrt{T \cdot 0.07}
= 3.59\ee{-9} \cdot 0.265 \sqrt{T}
= 9.5\ee{-10} \sqrt{T}.
\end{equation}

For $\SNR_1 = 1$: $T = (1/9.5\ee{-10})^2 = 1.1\ee{18}$~s $\approx 3.5\ee{10}$~years.
This is absurd for a single sensor.

For an array of $N$ sensors in a vertical borehole of depth $D$,
each separated by spacing $d$, with $N = D/d$:
\begin{equation}
\SNR_N = \sqrt{N} \cdot \SNR_1 \cdot \frac{D}{1\,\mathrm{m}}.
\end{equation}

The factor $D/1$~m arises because the total baseline $D$ produces
a larger gradient signal. More precisely, with $N$ sensors spanning
depth $D$, the gravity gradient can be extracted with uncertainty
reduced by $\sqrt{N}$ from averaging and signal enhanced by $D$:
\begin{equation}
\SNR_{\mathrm{array}} = \frac{D}{d} \cdot \frac{3.59\ee{-18} \cdot D}{10^{-9}/\sqrt{N}}
\cdot \sqrt{T \cdot \Delta f}.
\end{equation}

Wait --- this needs careful treatment. The measurable is the gravity
\emph{gradient}. An array of $N$ sensors at positions $z_i$ along a
vertical line measures $g(z_i)$. The DFD signal is a correction to
the gravity gradient:
\begin{equation}
\frac{\partial g_{\mathrm{DFD}}}{\partial z}
= \frac{\partial}{\partial z}\left(\frac{\azero}{g(z)} \cdot g(z)\right)
= \frac{\partial \azero}{\partial z} = 0.
\end{equation}

No --- $\azero$ is a constant. The DFD \emph{correction} to the
gradient is:
\begin{equation}
\frac{\partial}{\partial z}\left(\frac{g}{\mu(g/\azero)}\right)
= \frac{\partial g}{\partial z} \cdot \frac{1}{\mu}
+ g \cdot \frac{\partial}{\partial z}\left(\frac{1}{\mu}\right).
\end{equation}

In the Newtonian regime ($g/\azero \gg 1$, $\mu \approx 1 - \azero/g$):
\begin{equation}
\frac{1}{\mu} \approx 1 + \frac{\azero}{g},
\end{equation}
so the DFD correction to the gradient is:
\begin{equation}
\delta\left(\frac{\partial g}{\partial z}\right)
= \frac{\partial g}{\partial z} \cdot \frac{\azero}{g}
+ g \cdot \frac{\partial}{\partial z}\left(\frac{\azero}{g}\right)
= \frac{\partial g}{\partial z} \cdot \frac{\azero}{g}
- \azero \cdot \frac{1}{g}\frac{\partial g}{\partial z}
= 0.
\end{equation}

\begin{correction}[MEMS Gradient Signal Cancellation]
\label{cor:mems-gradient}
\corrected{The DFD $\mu$-function correction to the \emph{vertical
gravity gradient} $\partial g/\partial z$ cancels to leading order
in the Newtonian regime. This is because $\delta g = \azero$ is a
\textbf{constant offset} (independent of $z$ in the regime $g \gg \azero$),
so its gradient vanishes.}

The next-order correction is:
\begin{equation}
\delta\left(\frac{\partial g}{\partial z}\right)_{\mathrm{2nd}}
= \frac{\partial}{\partial z}\left(\frac{\azero^2}{g}\right)
= -\frac{\azero^2}{g^2}\frac{\partial g}{\partial z}
= \frac{2\azero^2}{R_E g}
= \frac{2 \times (1.12\ee{-10})^2}{6.371\ee{6} \times 9.81}
= 4.01\ee{-28}\,\mathrm{s}^{-2}.
\end{equation}

This is equivalent to $4.01\ee{-28}/9.81 \approx 4.1\ee{-29}$~g/m,
or $\sim 4\ee{-17}$~E\"otv\"os.

\caution{This second-order gradient signal is 20 orders below
current MEMS sensitivity. A MEMS array \textbf{cannot} detect the
DFD $\mu$-function gradient in the Newtonian regime. The W4i10
derivation of the MAGIS-100 gradient signal ($3.6\ee{-16}$~m/s$^2$)
is correct but refers to the \emph{absolute} DFD correction
accumulated over the baseline, not the differential gradient.
A gradiometer (which measures gradient, not absolute $g$) sees the
cancellation derived above.}
\end{correction}

\begin{finding}[MEMS Array Value for DFD]
\label{find:mems-value}
\honest{Chip-scale MEMS gravimeters at $0.1\;\mu\mathrm{Gal}/\sqrt{\mathrm{Hz}}$
represent a remarkable engineering achievement, but they \textbf{cannot}
detect the DFD $\mu$-function signal because:

\begin{enumerate}[nosep]
\item The DFD correction in the Newtonian regime is a constant offset
  $\delta g = \azero$, whose gradient is zero to leading order.
\item The second-order gradient is $\sim 4\ee{-17}$~E, undetectable.
\item Absolute $g$ measurements (not gradients) at $10^{-11}$ relative
  precision would detect the DFD signal, but systematic errors from
  local mass distribution uncertainties dominate at this level
  (same problem as MAGIS-100 DC offset, W4i10).
\end{enumerate}

The real DFD application of MEMS arrays is as \textbf{environmental
noise monitors}: distributed chip-scale gravimeters can map local mass
variations (groundwater, tides, atmospheric pressure) that constitute
Newtonian noise for other DFD experiments (atom interferometers, SRF
cavities). This is a support role, not a detection role.}
\end{finding}

%=============================================================================
\section{Finding 3: NASA QGG Pathfinder and Spaceborne DFD Testing}
\label{sec:qgg}
%=============================================================================

\subsection{NASA Quantum Gravity Gradiometer Pathfinder}

NASA funded (2024) the development of a Quantum Gravity Gradiometer
Pathfinder (QGGPf) instrument, with NET on-orbit testing 2030.
Key parameters (from EPJ Quantum Technology, March 2025):

\begin{itemize}[nosep]
\item \textbf{Configuration}: Single-axis cold-atom Rb gradiometer,
  cross-track orientation in LEO ($\sim$500~km altitude).
\item \textbf{Partners}: AOSense, Infleqtion, Vector Atomic, JPL.
\item \textbf{Science}: Validate long-term stability (1--90~min) of
  atomic gravity gradiometry in space against known Earth gravity field.
\item \textbf{Performance target}: Comparable or better resolution than
  GRACE-FO for mass change mapping.
\item \textbf{Goal}: Demonstrate feasibility of science-grade
  single-spacecraft gradiometer.
\end{itemize}

\subsection{DFD Signal at LEO Altitude}

At altitude $h = 500$~km above Earth's surface:
\begin{equation}
a_N(h) = g \left(\frac{R_E}{R_E + h}\right)^2
= 9.81 \times \left(\frac{6371}{6871}\right)^2
= 9.81 \times 0.859 = 8.43\,\mathrm{m/s}^2.
\end{equation}

The DFD correction:
\begin{equation}
\frac{\delta a}{a_N} = \frac{\azero}{a_N} = \frac{1.12\ee{-10}}{8.43}
= 1.33\ee{-11}.
\end{equation}

The ratio to the surface value:
\begin{equation}
\frac{(\delta a/a_N)_{500\,\mathrm{km}}}{(\delta a/a_N)_{\mathrm{surface}}}
= \left(\frac{6871}{6371}\right)^2 = 1.164.
\end{equation}

\caution{At 500~km altitude, the DFD correction is only $16.4\%$ larger
than at the surface. The altitude is far too low to access the MOND
transition region ($r_{\mathrm{cross}} \sim 13{,}000$--$22{,}000$~km).}

\subsection{DFD Gradient Signal on QGGPf Baseline}

The QGGPf measures the gravity gradient (difference between two atom
interferometers separated by baseline $d \sim 0.3$~m in cross-track).
From the correction in Section~\ref{sec:mems}, the DFD gradient
contribution at leading order cancels. The second-order DFD gradient
correction at 500~km altitude:
\begin{equation}
\delta T_{zz}^{\mathrm{DFD}} = \frac{2\azero^2}{R_E g(h)}
= \frac{2 \times (1.12\ee{-10})^2}{6.871\ee{6} \times 8.43}
= 4.33\ee{-28}\,\mathrm{s}^{-2}
\approx 4.3\ee{-19}\,\mathrm{E}.
\end{equation}

This is approximately $10^{7}$ orders of magnitude below the
projected QGGPf sensitivity ($\sim 10^{-12}$~E range).

\begin{finding}[QGGPf Cannot Detect DFD Directly]
\label{find:qgg-honest}
\honest{The NASA QGGPf at 500~km LEO cannot detect the DFD
$\mu$-function signal. The DFD gradient correction
($\sim 4\ee{-19}$~E) is 7 orders below the instrument sensitivity.}

\newresult{However, QGGPf provides three indirect benefits for DFD P6:}
\begin{enumerate}[nosep]
\item \textbf{Systematics characterization}: Long-term stability
  validation of atomic gradiometry in space directly informs the
  design of a future DFD-optimized space gradiometer at higher orbit.
\item \textbf{Technology maturation}: QGGPf matures cold-atom
  space hardware (atom sources, vacuum, laser systems) to TRL~6+,
  reducing costs for any subsequent DFD space mission.
\item \textbf{Altitude scaling validation}: Comparing QGGPf
  gravity measurements with ground truth validates the $r^{-2}$
  scaling of $a_N$ in the Earth's field to high precision, which
  constrains the DFD correction ($\propto r^2$) indirectly.
\end{enumerate}

The optimal DFD space mission is \emph{not} QGGPf but a dedicated
atom interferometer at GPS or GEO altitude, where $\delta a/a_N$
grows to $1.14\ee{-11} \times 17.4 = 2.0\ee{-10}$ (GPS) or
$\times 44.6 = 5.1\ee{-10}$ (GEO). At GEO, the DFD correction
to $g$ is $\delta a = 5.1\ee{-10} \times 0.22 = 1.1\ee{-10}$~m/s$^2$
$= \azero$ --- i.e., GEO altitude ($\sim 42{,}000$~km) is
approximately at the DFD crossover. An AEDGE-class mission at GEO
would be the definitive P6 experiment.
\end{finding}

%=============================================================================
\section{Finding 4: Gravitational Decoherence as DFD Null Test}
\label{sec:decoherence}
%=============================================================================

\subsection{Gravitational Decoherence in Optomechanical Systems}

A February 2026 paper (arXiv:2602.14841) develops a quantum estimation
framework for gravitational decoherence in optomechanical systems.
The key physics: a quantum superposition of mass $m$ at spatial
separation $\Delta x$ experiences decoherence from gravitational
interactions at a rate that depends on the gravitational model.

In GR, the dominant gravitational decoherence mechanism is the
Di\'osi-Penrose (DP) model:
\begin{equation}
\Gamma_{\mathrm{DP}} = \frac{G m^2}{\hbar} \cdot f(\Delta x / R),
\end{equation}
where $R$ is the object size and $f$ is a geometry-dependent factor.
For a sphere of radius $R$ with $\Delta x \ll R$:
$f \approx (\Delta x)^2 / (5 R^2)$.

\subsection{DFD-Specific Decoherence Channel}

In DFD, the psi-field introduces an additional decoherence channel.
A mass $m$ in a spatial superposition at positions $\mathbf{x}_1$ and
$\mathbf{x}_2$ experiences different psi-field values:
\begin{equation}
\Delta\psi = \nabla\psi \cdot (\mathbf{x}_2 - \mathbf{x}_1)
= \nabla\psi \cdot \Delta\mathbf{x}.
\end{equation}

The resulting differential phase accumulation rate:
\begin{equation}
\dot{\phi}_\psi = \frac{m c^2}{2\hbar} \Delta\psi
= \frac{m c^2}{2\hbar} \nabla\psi \cdot \Delta\mathbf{x}.
\end{equation}

In the DFD framework, $\nabla\psi$ is determined by the local
gravitational environment. At Earth's surface:
\begin{equation}
|\nabla\psi| = \frac{2g}{c^2} = \frac{2 \times 9.81}{(3\ee{8})^2}
= 2.18\ee{-16}\,\mathrm{m}^{-1}.
\end{equation}

The DFD decoherence rate (from differential phase randomization
due to psi-field fluctuations at frequency $\omega$) is:
\begin{equation}
\Gamma_\psi = \left(\frac{m c^2}{2\hbar}\right)^2
\cdot S_{\nabla\psi}(\omega) \cdot (\Delta x)^2,
\end{equation}
where $S_{\nabla\psi}(\omega)$ is the power spectral density of
psi-gradient fluctuations.

But in DFD, the psi-gradient is a \emph{classical} deterministic
field (not stochastic). Its fluctuations arise only from the
time-varying mass distribution. For a static laboratory:
$S_{\nabla\psi} \approx 0$ at frequencies above the seismic band.
The decoherence from the static gradient is:
\begin{equation}
\Gamma_\psi^{\mathrm{static}} = 0
\end{equation}
--- a static gradient produces a constant phase difference, not
decoherence.

For time-varying sources (e.g., seismic noise at $\omega \sim 1$~Hz):
\begin{equation}
S_{\nabla\psi}(\omega) = \left(\frac{2}{c^2}\right)^2 S_{\nabla g}(\omega)
= \frac{4}{c^4} S_{\nabla g}(\omega).
\end{equation}

With seismic gravity gradient noise
$S_{\nabla g}^{1/2} \sim 10^{-15}$~m/s$^2$/m/$\sqrt{\mathrm{Hz}}$
(typical for quiet sites):
\begin{equation}
\Gamma_\psi = \left(\frac{mc^2}{2\hbar}\right)^2
\cdot \frac{4}{c^4} \cdot (10^{-15})^2 \cdot (\Delta x)^2 \cdot \Delta f.
\end{equation}

For a 1.2~ng Si$_3$N$_4$ membrane ($m = 1.2\ee{-12}$~kg),
$\Delta x = 100$~nm, $\Delta f = 1$~Hz:
\begin{equation}
\Gamma_\psi = \left(\frac{1.2\ee{-12} \times (3\ee{8})^2}{2 \times 1.055\ee{-34}}\right)^2
\cdot \frac{4}{(3\ee{8})^4} \cdot 10^{-30} \cdot 10^{-14} \cdot 1.
\end{equation}

\begin{equation}
= \left(\frac{1.08\ee{5}}{2.11\ee{-34}}\right)^2 \cdot 4.94\ee{-78}
= (5.12\ee{38})^2 \cdot 4.94\ee{-78}
= 2.62\ee{77} \cdot 4.94\ee{-78}
= 1.3\,\mathrm{s}^{-1}.
\end{equation}

Wait --- this is suspiciously large. Let me recheck.

The issue is the enormous prefactor $(mc^2/2\hbar)^2$.
But this should not produce decoherence from the psi-gradient
directly; the decoherence comes from the \emph{fluctuating}
component of the gradient coupling into differential phase
\emph{noise}. The correct treatment: the psi-gradient couples
as a classical force to both arms of the superposition. In a
\emph{common-mode} geometry (same gravitational environment for
both superposition branches), the psi-gradient force is
\emph{identical} on both branches and produces no decoherence.

The decoherence arises only from the \emph{gradient of the gradient}
(tidal term) over the superposition separation $\Delta x$:
\begin{equation}
\delta F = m \cdot \frac{c^2}{2} \cdot \frac{\partial^2 \psi}{\partial x^2}
\cdot (\Delta x)^2,
\end{equation}
and the decoherence rate is:
\begin{equation}
\Gamma_\psi^{\mathrm{tidal}} = \frac{1}{\hbar^2}
\cdot \left(\frac{mc^2}{2}\right)^2
\cdot S_{\partial^2\psi/\partial x^2}(\omega)
\cdot (\Delta x)^4 \cdot \Delta f.
\end{equation}

With $\partial^2\psi/\partial x^2 = (2/c^2) \partial^2 g/\partial x^2
\sim (2/c^2)(g/R_E^2) \sim 4.8\ee{-30}\,\mathrm{m}^{-2}$ (static, no
fluctuation), the static tidal decoherence is again zero. The
fluctuating tidal component from seismic sources at lab scales
($\sim 10$~m) is:
\begin{equation}
S_{\partial^2 g/\partial x^2}^{1/2} \sim 10^{-15} / 10
= 10^{-16}\,\mathrm{s}^{-2}\mathrm{m}^{-1}\sqrtHz,
\end{equation}
giving $S_{\partial^2\psi/\partial x^2}^{1/2}
\sim (2/c^2) \times 10^{-16} = 2.2\ee{-33}\,\mathrm{m}^{-2}\sqrtHz$.

Then:
\begin{equation}
\Gamma_\psi^{\mathrm{tidal}} = \frac{(mc^2/2)^2}{\hbar^2}
\cdot (2.2\ee{-33})^2 \cdot (10^{-7})^4 \cdot 1
= (5.12\ee{38})^2 \cdot 4.8\ee{-66} \cdot 10^{-28}
\end{equation}
\begin{equation}
= 2.62\ee{77} \cdot 4.8\ee{-94}
= 1.3\ee{-16}\,\mathrm{s}^{-1}.
\end{equation}

But this is still the \emph{Newtonian} tidal decoherence, not a
DFD-specific effect. The DFD-specific contribution comes from the
$\mu$-function correction. The DFD correction to the tidal tensor is:
\begin{equation}
\delta T_{ij}^{\mathrm{DFD}} = \frac{\azero}{a_N} T_{ij}^N
\sim 10^{-11} T_{ij}^N.
\end{equation}

So the DFD-specific decoherence rate is:
\begin{equation}
\Gamma_\psi^{\mathrm{DFD}} = \left(\frac{\azero}{a_N}\right)^2 \Gamma_\psi^{\mathrm{tidal}}
= (10^{-11})^2 \times 1.3\ee{-16}
= 1.3\ee{-38}\,\mathrm{s}^{-1}.
\end{equation}

\begin{finding}[DFD Gravitational Decoherence: Null Prediction]
\label{find:decoherence}
\newresult{The DFD-specific contribution to gravitational decoherence
in a 1.2~ng membrane with 100~nm spatial superposition is
$\Gamma_\psi^{\mathrm{DFD}} \sim 10^{-38}$~s$^{-1}$.}

This is:
\begin{itemize}[nosep]
\item 32 orders below current experimental sensitivity
  ($\Gamma_{\mathrm{exp}} \sim 10^{-6}$~s$^{-1}$).
\item 22 orders below the standard Newtonian tidal decoherence
  ($\sim 10^{-16}$~s$^{-1}$).
\item Consistent with the DFD passive-dominance paradigm: the
  DFD correction is $(\azero/a_N)^2 \sim 10^{-22}$ times the
  Newtonian tidal effect.
\end{itemize}

\verdict{DFD predicts \textbf{zero anomalous gravitational decoherence}
at current and foreseeable experimental sensitivity. This is a robust
\emph{null prediction} that is consistent with all existing quantum
coherence measurements. Any observation of anomalous gravitational
decoherence at rates $> 10^{-16}$~s$^{-1}$ in the relevant mass/size
regime would be inconsistent with DFD (and with standard Newtonian
gravity), providing a potential falsification channel.}

\caution{This null prediction is also shared by GR and all other
theories that do not introduce new short-range gravitational
interactions. DFD's decoherence prediction is non-distinctive ---
it cannot be used to \emph{confirm} DFD, only to falsify alternative
theories that predict enhanced decoherence.}
\end{finding}

%=============================================================================
\section{Finding 5: Quantum Gradiometer Applications}
\label{sec:applications}
%=============================================================================

\subsection{Field-Deployed Quantum Gradiometers: State of the Art}

The Birmingham cold-atom gravity gradiometer achieved the following
in field conditions (Nature, 2022):

\begin{itemize}[nosep]
\item Statistical uncertainty: 20~E ($1\,\mathrm{E} = 10^{-9}$~s$^{-2}$).
\item Spatial resolution: 0.5~m.
\item Detected a 2~m buried tunnel at SNR~8, localizing its centre
  to within 20~cm.
\item Operational in non-laboratory environment (between two
  multi-story buildings, near a road).
\end{itemize}

Commercial developments include Delta.g (University of Birmingham
spin-off), targeting civil engineering (brownfield site inspection)
and telecommunications (underground utility mapping).

The Oxford--Imperial College hybrid cold-atom/MEMS gradiometer
programme targets GPS-denied navigation for maritime applications,
combining quantum precision with MEMS robustness.

\subsection{DFD $\mu$-Function Signal in Gradiometric Applications}

For a subsurface density anomaly (e.g., ore body, tunnel, submarine)
of mass $M$ at depth $d$ from the sensor, the Newtonian gravity
gradient signal is:
\begin{equation}
T_{zz}^N = \frac{2GM}{d^3}\,\mathrm{E},
\end{equation}
(for a point mass, approximately).

The DFD correction to this gradient signal in the Newtonian regime
($a_N \gg \azero$) is:
\begin{equation}
\delta T_{zz}^{\mathrm{DFD}} = \frac{\azero}{a_N} T_{zz}^N
+ \text{(second-order cancellation terms)}.
\end{equation}

At the Earth's surface, $a_N = g$ and $\azero/g = 1.14\ee{-11}$.
The DFD correction to a 100~E gradiometric signal from a buried
tunnel is:
\begin{equation}
\delta T_{zz}^{\mathrm{DFD}} \sim 1.14\ee{-11} \times 100\,\mathrm{E}
= 1.14\ee{-9}\,\mathrm{E}.
\end{equation}

Current sensitivity: 20~E. Required improvement: $\sim 2\ee{10}\times$.

\begin{finding}[Gradiometric DFD Detection: Impossible]
\label{find:gradiometer-apps}
\honest{The DFD $\mu$-function correction to gravity gradiometric
signals is $\sim 10^{-11}$ relative, requiring $10^{10}\times$
improvement over current quantum gradiometer sensitivity. This
correction is:

\begin{itemize}[nosep]
\item Undetectable for mineral exploration (ore bodies produce
  $\sim 1$--$100$~E signals; DFD correction $\sim 10^{-11}$--$10^{-9}$~E).
\item Undetectable for submarine detection (submarine mass
  $\sim 10^4$~tonnes at 300~m depth produces $\sim 0.1$~E;
  DFD correction $\sim 10^{-12}$~E).
\item Undetectable for tunnel detection (2~m tunnel produces
  $\sim 100$~E; DFD correction $\sim 10^{-9}$~E).
\end{itemize}

The DFD $\mu$-function simply adds a negligible fractional
correction ($\sim 10^{-11}$) to all Newtonian gravitational
signals at the Earth's surface. Practical gradiometry applications
are entirely Newtonian.}

\verdict{Quantum gradiometers serve DFD as \textbf{infrastructure},
not as detection instruments. Specifically:
\begin{enumerate}[nosep]
\item Environmental gravity noise mapping for DFD experiments
  (atom interferometers, SRF cavities).
\item Technology demonstration for cold-atom hardware that feeds
  into the MAGIS/AION detection pathway.
\item Gravity survey around MAGIS-100's NuMI shaft to characterize
  the local mass distribution systematic that limits the DC offset
  measurement (W4i10, Finding~1).
\end{enumerate}}
\end{finding}

\subsection{Submarine and Military Applications: Honest Assessment}

The quantum sensing community frequently cites submarine detection
and anti-stealth applications. For DFD:

\begin{itemize}[nosep]
\item DFD adds no detectable advantage to submarine detection via
  gravity gradiometry. The signal is purely Newtonian.
\item The DFD-specific scalar signals (psi-gradient, $\dpsiem$) are
  not produced by submarines (which are not significant EM sources
  relative to the solar/terrestrial gravitational psi-field).
\item The claim that DFD enables new military sensor modalities
  is \textbf{unsupported}. P6 is a fundamental physics experiment,
  not a defence technology.
\end{itemize}

%=============================================================================
\section{Cross-Reference Updates and Consistency Checks}
\label{sec:crossref}
%=============================================================================

\subsection{Updates to MASTER\_FINDINGS\_CORPUS}

\begin{table}[H]
\centering
\small
\caption{W4i13 P6 corrections and confirmations for the corpus.}
\begin{tabular}{@{}p{4cm}p{4cm}p{5cm}@{}}
\toprule
Item & Previous & W4i13 Update \\
\midrule
MAGIS-100 SNR & 540 (W4i10) & \confirmed{No change.} \\
MEMS threshold & $\Qpsi \geq 2.2\ee{8}$ (W4i11) & \confirmed{No change.} \\
Si$_3$N$_4$ $Q \geq 10^9$ & W4i10 & \confirmed{Additionally validated
  by optomechanical platform (arXiv:2601.02576).} \\
Advantage band & 0.04--0.11 Hz & \confirmed{No change. New high-frequency
  complement at 1--100 kHz via optomechanics (Finding 1).} \\
MAGIS-100 timeline & Commission $\sim$2028 & \confirmed{Atom sources
  delivery late 2026; installation through 2027; commissioning 2028.} \\
MEMS gravity gradient & $3.6\ee{-16}$~m/s$^2$ (W4i10) &
  \caution{This is the absolute correction over baseline, not a
  gradient. The DFD gradient signal \emph{cancels} to leading order
  for gradiometers (Finding 2).} \\
\bottomrule
\end{tabular}
\end{table}

\subsection{MIGA Status}

The MIGA (Matter-wave laser Interferometric Gravitation Antenna)
underground atom interferometer at LSBB Rustrel (France) has a
150~m horizontal baseline with three cold-atom sources. MIGA
targets the 0.1--1~Hz band for both GW detection pathfinding and
geophysical gravity gradient monitoring. As of March 2026, the
instrument appears to be in the installation/commissioning phase
(no published first-light results found). MIGA is complementary
to MAGIS-100 (vertical vs.\ horizontal baseline) and provides
an independent platform for DFD P6 scalar-mode searches.

\subsection{DECIGO/B-DECIGO Status}

B-DECIGO remains planned for launch in the 2030s. No significant
status changes since W4i10. The decihertz gap persists.
\confirmed{Finding 3 of W4i10 (decihertz gap through $\sim$2040)
remains valid.}

\subsection{Levitated Nanoparticle Force Sensing}

Levitated optomechanical systems have achieved force sensitivity
of $\sim 4.3$~zN/$\sqrt{\mathrm{Hz}}$ (zeptonewton regime).
These offer an alternative to membrane-based sensors for DFD
scalar detection, but are currently at lower TRL ($\sim$2) and
the force sensitivity maps to psi-gradient sensitivity of:
\begin{equation}
|\nabla\psi|_{\min}^{\mathrm{levitated}} = \frac{2 F_{\min}}{m c^2}
= \frac{2 \times 4.3\ee{-21}}{10^{-15} \times (3\ee{8})^2}
= 9.6\ee{-35}\,\mathrm{m}^{-1}\sqrtHz,
\end{equation}
where we assumed $m \sim 10^{-15}$~kg (nanosphere).

This is $\sim 10^4$ better than the membrane platform (Finding~1)
but applies at much lower frequencies ($\sim 10$~Hz, limited by
trap frequency). The levitated platform is a promising long-term
alternative but not yet competitive with atom interferometry for
DFD-relevant frequency bands.

%=============================================================================
\section{Inconsistencies Flagged}
\label{sec:inconsistencies}
%=============================================================================

\begin{enumerate}
\item \textbf{MAGIS-100 ``gradient signal'' terminology (MINOR)}:
  The W4i10 derivation of $\delta(\delta a) = 3.6\ee{-16}$~m/s$^2$
  computes the \emph{differential DFD correction} between top and
  bottom of the 100~m baseline. This is correctly described as a
  ``gradient signal'' in the context of an atom interferometer
  (which measures differential phase, hence differential acceleration).
  However, it should be clarified that this is \emph{not} the same as
  the gravity gradient $\partial g/\partial z$ measured by a
  gradiometer. A gradiometer measures the derivative of $g$; the DFD
  correction to this derivative cancels to leading order (Finding~2).
  The atom interferometer measures the \emph{absolute} acceleration
  difference, which does not cancel.

\item \textbf{GEO altitude crossover (NEW, MINOR)}:
  Finding~3 derives that GEO altitude ($\sim 42{,}000$~km from
  Earth's centre) corresponds to $\delta a/a_N \sim 5\ee{-10}$
  and $\delta a \sim \azero$. This is \emph{approximately} the
  crossover, but the exact crossover altitude (where $a_N = \azero$)
  is at $r = \sqrt{GM_E/\azero} = \sqrt{3.986\ee{14}/1.12\ee{-10}}
  = \sqrt{3.56\ee{24}} = 1.89\ee{12}$~m $= 1.89\ee{6}$~km.
  This is much larger than GEO ($42{,}164$~km). At GEO, $a_N = 0.22$~m/s$^2$,
  giving $a_N/\azero = 2\ee{9}$, still deep in the Newtonian regime.

  \corrected{The W4i10 ``crossover altitude'' of 13,000--22,000~km
  (from W3i3) refers to $r_{\mathrm{cross}} \sim 13{,}000$--$22{,}000$~km
  \emph{from Earth's centre}. Let me verify: $r = \sqrt{GM_E/\azero}
  = \sqrt{3.986\ee{14}/1.12\ee{-10}} = 1.89\ee{12}$~m $= 1.89\ee{6}$~km
  $= 1.89$~million~km. This is clearly NOT 13,000--22,000~km.}

  \textbf{CRITICAL CHECK}: The W3i3 crossover altitude must have
  used a different definition. Reviewing the W4i10 Finding~5:
  ``$r_{\mathrm{cross}} \approx 13{,}000$--$22{,}000$~km
  (inherited W3i3).'' This requires verification.

  For $a_N = \azero$ around Earth:
  $GM_E/r^2 = \azero$, so $r = \sqrt{GM_E/\azero}
  = \sqrt{3.986\ee{14}/1.12\ee{-10}} = \sqrt{3.56\ee{24}}
  = 5.96\ee{7}$~km $\approx 60{,}000$~km? Let me recompute:
  $\sqrt{3.56\ee{24}} = \sqrt{3.56} \times 10^{12} = 1.887\ee{12}$~m
  $= 1.887\ee{9}$~m ... no.

  $\sqrt{3.56\ee{24}\,\mathrm{m}^3/\mathrm{s}^2 \;/\;
  (1.12\ee{-10}\,\mathrm{m/s}^2)}$ ...
  Units: $\sqrt{\mathrm{m}^3\mathrm{s}^{-2}/(\mathrm{m\,s}^{-2})}
  = \sqrt{\mathrm{m}^2} = \mathrm{m}$.

  $r = \sqrt{3.986\ee{14}/1.12\ee{-10}} = \sqrt{3.56\ee{24}}$~m.

  $\sqrt{3.56\ee{24}} = \sqrt{3.56} \times 10^{12}
  = 1.887 \times 10^{12}$~m $= 1.887\ee{9}$~km.

  This is $\sim 1.9$ \emph{billion} km, approximately 12.6~AU.
  The MOND crossover for Earth's gravity is at 12.6~AU from Earth.
  This makes physical sense --- MOND effects become important only
  at astronomical distances where $a_N \sim \azero$.

  \begin{correction}[Crossover Altitude for Earth's Field]
  \label{cor:crossover}
  \corrected{The MOND crossover for Earth's gravitational field
  (where $a_N(r) = \azero$) is at $r = 1.89\ee{9}$~km $\approx 12.6$~AU,
  not 13,000--22,000~km. The W3i3 value of 13,000--22,000~km is
  \textbf{incorrect by a factor of $\sim 10^5$}.}

  \caution{This error propagated through W3i3, W4i10 Findings~2
  and 5, and potentially other documents. The W4i10 prediction
  of ``crossover altitude accessible by satellite atom interferometer''
  is invalid --- the crossover is at interplanetary distances, not
  Earth orbit.}

  The correct interpretation: the W3i3 value may have referred to
  the altitude where the DFD correction becomes detectable at some
  assumed sensitivity, not where $a_N = \azero$. However, the text
  explicitly states ``$a_N = \azero$, $r_{\mathrm{cross}} \approx
  13{,}000$--$22{,}000$~km.'' This is a mathematical error:
  $GM_E / (2\ee{7}\,\mathrm{m})^2 = 3.986\ee{14}/(4\ee{14})
  \approx 1$~m/s$^2$, not $1.12\ee{-10}$~m/s$^2$.

  At 20,000~km from Earth's centre: $a_N = GM_E/r^2
  = 3.986\ee{14}/(2\ee{7})^2 = 0.997$~m/s$^2$.

  \verdict{The crossover altitude in the W3i3/W4i10 sense requires
  $a_N = \azero$. With Earth's field alone, this is at 12.6~AU.
  The 13,000--22,000~km figure is wrong.
  This is a MAJOR inherited error. Must be flagged for corpus correction.}
  \end{correction}

\item \textbf{No other inconsistencies found} in the W4i10 P6 document
  beyond the crossover altitude error.
\end{enumerate}

%=============================================================================
\section{Summary of New Territory Explored}
\label{sec:summary}
%=============================================================================

\begin{table}[H]
\centering
\small
\caption{New sensor modalities assessed for DFD P6 (W4i13).}
\begin{tabular}{@{}p{3.5cm}p{2cm}p{2cm}p{5cm}@{}}
\toprule
Modality & DFD Signal? & Detectable? & Assessment \\
\midrule
Si$_3$N$_4$ optomechanical membrane & Yes (scalar) & Not yet &
  Best high-frequency platform; $10^4\times$ short of astrophysical \\
MEMS chip-scale array & Gradient cancels & No &
  Useful as environmental noise monitor only \\
NASA QGGPf (LEO) & Yes (tiny) & No &
  $10^7\times$ short; value as technology maturation \\
Levitated nanoparticle & Yes (scalar) & Not yet &
  $10^4\times$ better psi-gradient sensitivity than membrane;
  lower TRL \\
Gravitational decoherence & Null prediction & N/A &
  DFD predicts zero anomalous decoherence; non-distinctive \\
Mineral/submarine gradiometry & $10^{-11}$ relative & No &
  Newtonian gravity dominates completely \\
\bottomrule
\end{tabular}
\end{table}

%=============================================================================
\section*{Recommendations for W4i14+}
%=============================================================================

\begin{enumerate}
\item \textbf{URGENT}: Verify and correct the inherited
  ``crossover altitude 13,000--22,000~km'' claim from W3i3.
  If this is an error (as derived here), it propagates through
  multiple documents and affects the AEDGE/satellite mission case.

\item \textbf{Optomechanical convergence}: Engage with the
  Rousso et al.\ (arXiv:2601.02576) platform as a potential
  host for DFD scalar-mode searches. The Si$_3$N$_4$ hardware
  overlap with the P15 MEMS programme is a significant synergy.

\item \textbf{MEMS environmental monitoring}: The 0.1~$\mu$Gal
  chip-scale MEMS gravimeter should be evaluated as a Newtonian
  noise monitor for the MAGIS-100 site at Fermilab. Cost is minimal
  ($\sim$\$10K per sensor).

\item \textbf{MIGA engagement}: Monitor MIGA first-light results
  (expected $\sim$2026--2027) for DFD P6 horizontal-baseline
  complementarity.

\item \textbf{Levitated optomechanics}: Track zeptonewton force
  sensing developments for long-term DFD scalar-mode detection
  below 100~Hz.
\end{enumerate}

\end{document}
